\section{Mobile Delivery, Responsive Qualification and Sovereign Transport Hardening}
\label{sec:pilot-sms-viewport-call-wave-hardening}

This implementation pass closes five independently executable tasks from the unified
mobile programme: representative responsive viewport qualification
(\texttt{MOB-0216}), an official SMS adapter (\texttt{MOB-0702}), adverse-network
and authority tests for calls and groups (\texttt{MOB-1208}), reuse of the existing
Wave framing and queue foundations (\texttt{MOB-1302}), and radio-bridge hardening
(\texttt{MOB-1304}).  Provider accounts, independently audited media encryption,
direct-radio hardware and internet-disconnected field delivery remain explicit
external gates.

\subsection{Official SMS delivery boundary}

Pilot now has a provider adapter for the Twilio Messages API.  The adapter is
selected by a persona-bound service configuration on the mother brain; credentials
are never sent to the television, phone interface or recipient.  SMS and WhatsApp
are separate configured channels even when both use the same provider API.  A
contact's SMS number is stored and resolved independently from its email, Pilot and
WhatsApp routes.

Before submission, Pilot validates an exact E.164 destination, a bounded message
body and the identity whose provider account will be used.  The approved draft
identifier becomes the provider idempotency key.  A successful API response must
contain a valid provider account identifier, message identifier and recognized
provider state.  Pilot retains only a narrow receipt and never treats API acceptance
as carrier delivery, human reading, task acceptance or consent.

Where a public HTTPS status-callback endpoint is configured, incoming callback
parameters are authenticated with the provider's HMAC signature over the exact
callback URL and sorted form values.  Invalid signatures fail closed.  Live SMS
qualification still requires the customer's authorized provider account, a real
recipient and callback receipts.  The WhatsApp route remains unavailable until an
officially approved sender and account are supplied; browser automation and
unofficial session extraction are not substitutes.

\subsection{Representative responsive mobile qualification}

The unified mobile surface was exercised in an automated real browser at three
representative CSS viewport sizes: 320 by 568 pixels, 360 by 800 pixels and 430 by
932 pixels.  These cover a small iPhone-class display, a common Android-class
display and a large iPhone-class display.  At every size the document width equalled
the viewport width, no horizontal overflow occurred, the application landmark and
primary navigation remained available, and every visible interactive target was at
least 44 by 44 CSS pixels.

The shared style contract now enforces the minimum target on mode controls, section
headers, cards, the composer and icon/avatar buttons.  This closes responsive layout
qualification, but does not replace later physical-device testing of VoiceOver,
TalkBack, guardian/child journeys, camera permissions, biometrics and background
lifecycle behavior.

\subsection{Adverse-network and compromised-participant behavior}

The call authority was tested with missing and out-of-order transport records,
simulated loss and jitter, reconnect epochs, participant removal and forged-device
attempts.  A call cannot connect until both parties have registered possession-signed
transport records for the current epoch.  Delayed legitimate evidence can recover a
partial negotiation; stale evidence cannot cross a reconnect boundary because every
reconnect advances the epoch and requires fresh keys, fingerprints and nonces.

Group messages now carry an exact canonical payload signed by the sending device.
The mother verifies current membership, the current group-key epoch, payload bounds
and a single-use nonce.  A removed member cannot post under the new epoch, a replayed
message is rejected, and possession of public identifiers alone cannot forge another
participant.  These tests qualify the signaling and authorization state machine.
They do not close \texttt{MOB-1201}: the native WebRTC media construction still
requires independent end-to-end-encryption review.

\subsection{Reuse of the Wave foundation}

The optional Wave path now exposes an inspectable foundation receipt describing the
existing reusable machinery: bounded indexed fragments, per-fragment and complete
payload SHA-256 integrity, Ed25519-authenticated bundles, priority/TTL/retry and
idempotent queueing, duplicate suppression, and store-and-forward persistence.
Local reconstruction remains explicitly different from physical delivery.  The
ordinary Pilot product continues to hide Wave, PHO and wallet functions unless
supported hardware and policy deliberately enable them.

The legacy GlyphNet radio node retains its genuine reconnecting serial driver and
authenticated WebSocket bridge driver.  A real link automatically disables the
development mock.  Development injection routes now return not-found unless an
explicit development flag enables them.  The service has no default bridge secret,
does not accept bridge credentials in WebSocket query strings and requires a header
token plus a versioned timestamp, single-use nonce and HMAC for every radio WebSocket
connection.  Nonces are bounded, expire with the authentication window and are
remembered to prevent replay.

The radio node was migrated to Express 5.2.1, recompiles successfully and reports no
known production dependency vulnerabilities in the package audit.  This closes the
software bridge-hardening task without claiming a radio product exists.  Genuine BLE
and Wi-Fi Direct drivers, ESP32 or Raspberry Pi qualification, regional frequency
compliance, and two-mother text/push-to-talk delivery with the internet physically
disabled remain \texttt{MOB-1307--MOB-1310}.

\subsection{Acceptance boundary}

Automated evidence covers exact SMS routing and callback authentication, contact
channel separation, responsive target sizing, signed adverse-network call and group
behavior, Wave framing/queue integrity, and source-level radio security invariants.
The complete regression contains 588 passing AION Fabric tests.  The radio-node
TypeScript build passes and its production dependency audit reports zero known
vulnerabilities.  This append-only section also compiles independently.
No message was sent, no external account was connected, no physical radio delivery
was attempted, and the live household television runtime was not restarted during
this work.
